\documentclass[conference]{IEEEtran}

\usepackage{amsmath,amssymb,bm}
\usepackage{booktabs}
\usepackage{graphicx}
\usepackage{cite}
\usepackage{xcolor}
\usepackage{url}
\usepackage{multirow}
\usepackage[caption=false,font=footnotesize]{subfig}

\newcommand{\todo}[1]{\textcolor{red!70!black}{\textbf{[TODO:} #1\textbf{]}}}
\newcommand{\rev}[1]{\textcolor{red}{#1}}
\newcommand{\calD}{\mathcal{D}}
\newcommand{\feasSet}{\mathcal{F}}
\newcommand{\AURCref}{\mathrm{AURC}_{\mathrm{ref}}}
\newcommand{\Nfive}{N_{5\%}}
\newcommand{\distdyn}{d^{\mathrm{dyn}}}
\newcommand{\placeholderfigurewide}[1]{%
    \fbox{\parbox[c][4.5cm][c]{0.96\textwidth}{\centering\small
    \textbf{FIGURE PLACEHOLDER}\\[1.5mm]#1}}}
\newcommand{\figurepdf}[2]{%
    \IfFileExists{#2}{%
        \includegraphics[width=#1]{#2}%
    }{%
        \fbox{\parbox[c][3.4cm][c]{#1}{\centering\small
        \textbf{PDF FIGURE PENDING}\\[1.5mm]\texttt{#2}}}%
    }%
}

\title{Fleet-Informed Bayesian Optimization for Data-Efficient Controller Calibration}

\author{\IEEEauthorblockN{Jakob Weber, Markus Gurtner, Adrian Trachte, Andreas Kugi}
\IEEEauthorblockA{\textit{Draft author list -- affiliations to be finalized}}}

\begin{document}
\maketitle

\begin{abstract}
\rev{Controller calibration is often repeated for every new system variant, although previously calibrated fleet members already contain information about promising controller parameters. This paper investigates whether such fleet knowledge can reduce the number of new closed-loop experiments required to calibrate another system. We consider linear-quadratic integral (LQI) yaw-rate control for a heterogeneous vehicle fleet and use Bayesian optimization (BO) to tune three controller-weight parameters. Historical scalar calibration data from previously calibrated systems initialize the Gaussian-process surrogate of a new system, while every proposed parameter vector is converted into an LQI gain using that system's own identified model. Across 120 new-system calibration problems from 20 held-out fleet realizations, fleet-informed BO reduces the budget-restricted mean number of experiments required to reach within 5\% of an independently scored offline reference from 6.72 to 3.08, a reduction of 54\%. All warm-started runs reach the target within the 12-experiment budget, compared with 79.2\% under independent BO. Dynamics-aware weighting of the historical data provides no measurable improvement over unweighted pooling. A separate 128-point landscape study shows that 99.4\% of evaluated system pairs retain Spearman rank correlation above 0.95 on their jointly feasible calibration points despite substantial dynamics heterogeneity, although a modest residual relation with dynamics similarity remains. The historical data require experiments on earlier fleet members, but this cost is incurred once and can be reused when later systems are commissioned.}
\end{abstract}

\begin{IEEEkeywords}
Bayesian optimization, controller calibration, transfer learning, LQI, heterogeneous systems, vehicle control
\end{IEEEkeywords}

\section{Introduction}
Controller calibration is an important but expensive part of commissioning modern dynamical systems. Even when the controller structure follows from model-based design, parameters such as linear-quadratic regulator or integral-control weights determine the final compromise between tracking performance, control effort, and robustness. In practice, these parameters are often refined through repeated simulation or closed-loop testing. Bayesian optimization (BO) is attractive for this task because it can optimize low-dimensional controller parameters from a limited number of expensive evaluations~\cite{frazier2018tutorial,marco2016lqr,berkenkamp2016safe,schillinger2017safe,vonrohr2024local}.

A fleet setting changes the calibration problem fundamentally. Vehicles or machines are commissioned one after another, so previously calibrated clients already contain closed-loop evaluations collected during earlier commissioning. Repeating a cold-start optimization for every new client discards this accumulated information. Reusing historical records does not eliminate the cost of generating them. Instead, that cost is incurred on earlier clients and can be amortized across later fleet members, while the new client requires only the experiments that must be executed on its own plant.

\rev{The contribution of this paper is therefore not a new BO acquisition rule. It is a control-oriented transfer formulation and evaluation of how previously collected controller-calibration experiments can be reused across heterogeneous dynamical systems while every new system still synthesizes its own feedback controller from a local model.}

\subsection{Related work and positioning}
\rev{BO has been used for automatic controller tuning in LQR, PI, and safety-constrained calibration problems~\cite{marco2016lqr,berkenkamp2016safe,schillinger2017safe,vonrohr2024local}. Transfer-learning BO reuses observations or surrogate information from related optimization tasks~\cite{joy2019transfer,tighineanu2022transfer,perrone2018scalable}, while collaborative and federated BO share information across multiple agents~\cite{dai2020fts,dai2021dpfts,chen2025collaborative}. These methods establish that source information can accelerate black-box optimization. They do not, however, answer the control-specific question studied here: whether closed-loop calibration data collected on physically different plants remain useful when the feedback gain is recomputed from each plant's own model. Our study focuses on this transferability question, its incremental commissioning benefit, and the relation between plant heterogeneity and calibration-landscape heterogeneity.}

The answer is not obvious. Heterogeneous vehicle dynamics can change closed-loop behavior substantially, and uncertainty-aware dynamics estimates contain useful structure for clustered system identification and controller design~\cite{weber2026ifac}. This motivates dynamics-aware transfer, where observations from dynamically similar clients are trusted more strongly. At the same time, the controller-calibration problem may be more invariant than the plant itself. Different plant models can induce similar rankings of the same low-dimensional LQI weight configurations. In that case, broad reuse of fleet experience may be sufficient and explicit source selection may add little.

This paper studies that distinction for an existing fleet and a newly arriving client. Previously calibrated clients, which we call \emph{source clients}, provide scalar records containing an evaluated controller-weight vector and its normalized closed-loop cost. The new client begins with two local initialization experiments. It then fits a Gaussian-process (GP) surrogate using its own measurements together with the source records and sequentially evaluates new LQI calibrations. The historical source records were collected before the new client is commissioned. They therefore require no new source-client experiments during this process. Only the new client's own closed-loop evaluations are counted when we report commissioning effort. We transfer calibration experience in the common weight space, not feedback gains. For every proposed weight vector, the new client computes its own LQI gain from its identified model before evaluating the controller.

The final evaluation protocol, including the compared methods, local experiment budget, primary metric, and statistical tests, was fixed before the 20 final fleet realizations were examined. The contributions are:
\begin{itemize}
    \item We formulate \emph{fleet-informed controller calibration} as a transfer problem in which historical calibration observations from previously commissioned clients are reused to reduce the number of new closed-loop experiments required by another client.
    \item \rev{In a held-out evaluation on a heterogeneous nonlinear vehicle fleet, fleet-informed BO reduces the budget-restricted mean number of local evaluations needed to reach within 5\% of an independently scored offline reference from 6.72 to 3.08, while all warm-started systems reach the criterion within the 12-evaluation budget.}
    \item \rev{We characterize why broad transfer is possible. Although the identified vehicle dynamics differ substantially, 99.4\% of pairwise LQI calibration landscapes remain rank-correlated above 0.95 on jointly feasible calibration points. Dynamics similarity retains a modest residual relation with landscape similarity, but dynamics-aware weighting still provides no measurable improvement over simple unweighted pooling on the held-out fleets.}
\end{itemize}

Figure~\ref{fig:concept} summarizes the commissioning workflow.

\begin{figure*}[t]
    \centering
    \includegraphics[width=0.96\textwidth]{figures/concept.pdf}
    \caption{Fleet-informed controller calibration. Historical calibration records from previously commissioned source clients initialize the Bayesian-optimization surrogate of a newly arriving system. Only new closed-loop evaluations performed on that system contribute to its commissioning cost. Each feedback gain is synthesized from the new system's own identified model.}
    \label{fig:concept}
\end{figure*}

\section{Problem Formulation}

\subsection{Heterogeneous systems and local controller design}
Consider systems $i\in\mathcal C$ with true nonlinear plants
\begin{align}
 \mathbf{x}_{i,k+1}&=\mathbf{f}_i(\mathbf{x}_{i,k},u_{i,k},\mathbf{w}_{i,k}),\\
 \mathbf{y}_{i,k}&=\mathbf{h}_i(\mathbf{x}_{i,k},\mathbf{v}_{i,k}),
 \label{eq:true_plant_generic}
\end{align}
where $\mathbf{f}_i$ differs across systems. Each system identifies a compact linear control-design model
\begin{equation}
 \mathbf{x}_{k+1}=\mathbf{A}(\boldsymbol{\theta}_i)\mathbf{x}_k+\mathbf{B}(\boldsymbol{\theta}_i)u_k.
 \label{eq:identified_model}
\end{equation}
Local recursive least-squares identification yields the parameter estimate $\hat{\boldsymbol{\theta}}_i$ and the associated parameter covariance $\hat{\boldsymbol{\Sigma}}_i$. In the vehicle study, the state is $\mathbf{x}=[\beta,r]^\top$, where $\beta$ is vehicle sideslip and $r$ is yaw rate. \rev{The measured output used for LQI tracking is $r$, hence $\hat{\mathbf C}_i=[0\;1]$ for the identified two-state model.}

For yaw-rate tracking, the state is augmented with the integral tracking error
\begin{equation}
 z_{k+1}=z_k+T_s\left(r_k^{\rm ref}-r_k\right).
 \label{eq:integral_state}
\end{equation}
With $\mathbf{x}_k^a=[\mathbf{x}_k^\top,z_k]^\top$, the LQI penalty is parameterized by the three-dimensional vector $\boldsymbol{\xi}\in\Xi\subset\mathbb R^3$. \rev{The three entries $\xi_1$, $\xi_2$, and $\xi_3$ parameterize the penalties on sideslip $\beta$, yaw rate $r$, and integral tracking error $z$, respectively.}

The augmented-state channels have different units and characteristic magnitudes. We use the fixed scale vector
\begin{equation}
 \mathbf{s}=[0.053229,\,0.575573,\,0.139862]^\top
 \label{eq:state_scales}
\end{equation}
and define
\begin{align}
 \mathbf{Q}(\boldsymbol{\xi})&=\operatorname{diag}\left(\frac{10^{\xi_1}}{s_1^2},\frac{10^{\xi_2}}{s_2^2},\frac{10^{\xi_3}}{s_3^2}\right),\\
 \mathbf{R}&=1.
 \label{eq:lqi_weights}
\end{align}
The scaling places the three state channels on comparable physical ranges, so $\boldsymbol{\xi}$ mainly controls relative penalty strength rather than compensating for different units. \rev{The numerical scales were fixed before the final BO study from representative nominal closed-loop trajectories and were not tuned separately for individual vehicle families.}

For system $i$, each candidate $\boldsymbol{\xi}$ is converted into a system-specific gain
\begin{equation}
 \mathbf{K}_i(\boldsymbol{\xi})=\operatorname{LQI}\!\left(\hat{\mathbf{A}}_i,\hat{\mathbf{B}}_i,\hat{\mathbf{C}}_i,\mathbf{Q}(\boldsymbol{\xi}),\mathbf{R}\right).
 \label{eq:lqi_gain}
\end{equation}
The same calibration vector can therefore produce a different feedback gain for every system. The transferable quantity is the location and observed performance of a calibration point in the common $\boldsymbol{\xi}$-space.

For calibration episode $\omega$, applying $\mathbf{K}_i(\boldsymbol{\xi})$ to the true plant yields the yaw-rate tracking cost $J_i(\boldsymbol{\xi},\omega)$. The system objective is
\begin{equation}
 f_i(\boldsymbol{\xi})=\mathbb E_{\omega}\!\left[J_i(\boldsymbol{\xi},\omega)\right].
 \label{eq:client_objective}
\end{equation}
Only feasible controller candidates are considered valid calibration outcomes. The feasibility conditions are defined in Sec.~\ref{subsec:vehicle_task}.

\subsection{Fleet-informed calibration objective}
A previously calibrated source system $j$ provides records
\begin{equation}
 \calD_j^{\rm hist}=\left\{\left(\boldsymbol{\xi}_{j,m},\tilde y_{j,m}\right)\right\}_{m=1}^{M_j},
 \label{eq:historical_data}
\end{equation}
where
\begin{equation}
 \tilde y_{j,m}=\frac{J_j(\boldsymbol{\xi}_{j,m})}{J_j^{\rm base}}
 \label{eq:historical_normalization}
\end{equation}
is the cost normalized by the fixed baseline controller of source system $j$. This gives historical observations from systems with different absolute tracking-error levels a common scale.

For a newly arriving system $i$, all historical records belonging to that same system are excluded from the source pool. This held-out construction emulates a genuinely new system and prevents leakage from a completed calibration history of the plant that is currently being calibrated.

\rev{We distinguish the incremental experimental cost $N_i^{\rm new}$, which counts only closed-loop evaluations newly executed on system $i$, from the historical pool size $N_i^{\rm hist}$, which counts source observations available before system $i$ arrives. In the final full-pool study, $N_i^{\rm hist}=29\times25=725$. The paper's sample-efficiency claim concerns a reduction in $N_i^{\rm new}$ rather than a reduction in the total number of experiments ever performed across the fleet.}

\subsection{Offline reference and threshold metric}
\rev{For reporting, each held-out system receives a high-accuracy offline reference that is computed independently of the BO decision process. The reference search first evaluates 1024 Sobol points over $\Xi=[-3,3]^3$ on the system's calibration-episode bank. The eight best feasible Sobol candidates initialize bounded Powell refinements. The best resulting calibration vector $\boldsymbol{\xi}_i^{\rm ref}$ is then scored on an independent reference-test bank containing 20 stochastic episodes per maneuver, yielding $J_i^{\rm ref}$. The reference is used only for reporting and never enters GP fitting, acquisition, normalization, or infeasibility handling. A nested convergence audit on one target from each vehicle family showed that the best 1024-point Sobol value differed from the corresponding 4096-point value by only 0.018\% on average.}

\rev{To compare the BO trajectories and the reference on the same independent evaluation distribution, every incumbent stored along each completed BO trajectory is rescored on the same reference-test bank. Let $J_{i,\mathrm{test}}^{\mathrm{best}}(n)$ denote the best independently rescored incumbent available after $n$ new-system experiments. The latent 5\% hitting time is}
\begin{equation}
 T_{i,5\%}=\min\left\{n:J_{i,\mathrm{test}}^{\mathrm{best}}(n)\le1.05J_i^{\rm ref}\right\}.
 \label{eq:n5}
\end{equation}
\rev{All runs are followed to the same local budget $B=12$. If the threshold is not reached by $B$, the hitting time is right-censored at that budget. For budget-restricted mean summaries we use $\min(T_{i,5\%},B)$, denoted by $N_{i,5\%}$ in the remainder. Under this common Type-I censoring this is the restricted mean time to threshold over the available budget. Cumulative success curves treat censored runs as unresolved rather than as successes at $B$.}

\section{Fleet-Informed Bayesian Optimization}

\subsection{Warm-started GP surrogate}
The normalization in~\eqref{eq:historical_normalization} is also applied to every observation collected by the new system. The baseline controller is fixed before BO begins and is independent of the offline reference used for reporting. Infeasible evaluations receive the same fixed penalty in normalized units.

At local evaluation $n$, system $i$ fits a GP surrogate to
\begin{equation}
 \calD_i^{(n)}=\calD_i^{\rm local}(n)\cup\bigcup_{j\ne i}\calD_j^{\rm hist}.
 \label{eq:pooled_dataset}
\end{equation}
After initialization, the next calibration vector is selected with expected improvement~\cite{frazier2018tutorial},
\begin{equation}
 \boldsymbol{\xi}_{i,n+1}\in\arg\max_{\boldsymbol{\xi}\in\Xi}\operatorname{EI}_i\left(\boldsymbol{\xi}\mid\calD_i^{(n)}\right).
 \label{eq:ei}
\end{equation}
The new system then synthesizes $\mathbf{K}_i(\boldsymbol{\xi}_{i,n+1})$ and evaluates the controller on its own nonlinear plant. The source records remain fixed and require no additional source-system rollouts during this process.

\rev{The search region is the common box $\Xi=[-3,3]^3$. The local initialization uses two points of a Sobol low-discrepancy sequence. The source histories were generated with the same calibration coordinates and objective normalization, which allows them to be inserted directly into the new system's surrogate.}

\rev{For reproducibility, Table~\ref{tab:bo_details} summarizes the BO implementation. Four numerical entries that depend on the original experiment code are deliberately left as explicit verification items in this draft rather than reconstructed from memory. They should be replaced from the experiment configuration before submission.}

\begin{table}[t]
\centering
\caption{BO implementation used in the final study. Red entries require direct source-code verification before submission.}
\label{tab:bo_details}
\scriptsize
\begin{tabular}{@{}ll@{}}
\toprule
Quantity & Setting\\
\midrule
Search box & $[-3,3]^3$\\
Local initialization & 2 Sobol points\\
Acquisition & Expected improvement, minimization\\
Output scaling & $J_i/J_i^{\rm base}$\\
Source weighting floor & $0.02$\\
GP kernel & \rev{\textbf{verify exact kernel}}\\
GP base noise $\alpha_0$ & \rev{\textbf{verify exact value}}\\
EI maximization & \rev{\textbf{verify candidate/restart rule}}\\
Infeasible penalty & \rev{\textbf{verify normalized value}}\\
Baseline $\boldsymbol{\xi}^{\rm base}$ & \rev{\textbf{verify exact vector}}\\
\bottomrule
\end{tabular}
\end{table}

We compare three methods. Independent BO uses only observations collected on the new system. Global fleet warm start combines all available source records with equal weight. Dynamics-weighted warm start uses the same records but attenuates observations from dynamically dissimilar source systems. The latter is included as a mechanistic ablation.

\subsection{Dynamics-weighted ablation}
Because local identification provides both $\hat{\boldsymbol{\theta}}_i$ and $\hat{\boldsymbol{\Sigma}}_i$, we quantify pairwise dynamics similarity without using the latent family labels. \rev{The dynamics-weighted method therefore has access only to identified-model information and never receives the nominal, payload, or tire-degraded family identity.}
For systems $i$ and $j$,
\begin{equation}
(\distdyn_{ij})^2=\left(\hat{\boldsymbol{\theta}}_i-\hat{\boldsymbol{\theta}}_j\right)^\top\left(\hat{\boldsymbol{\Sigma}}_i+\hat{\boldsymbol{\Sigma}}_j+\epsilon\mathbf{I}\right)^{-1}\left(\hat{\boldsymbol{\theta}}_i-\hat{\boldsymbol{\theta}}_j\right),
\label{eq:dyn_distance}
\end{equation}
and
\begin{equation}
 s_{ij}=\exp\left(-\frac{(\distdyn_{ij})^2}{2\ell_d^2}\right),
 \label{eq:dyn_similarity}
\end{equation}
where $\ell_d$ is the median pairwise-distance bandwidth. Lower source similarity is converted into larger GP observation noise,
\begin{equation}
 \alpha_{ij}=\frac{\alpha_0}{\max(s_{ij},0.02)}.
 \label{eq:weighted_noise}
\end{equation}
\rev{The regularization $\epsilon$ in~\eqref{eq:dyn_distance} and the base noise $\alpha_0$ in~\eqref{eq:weighted_noise} are fixed across all systems and methods. Their exact numerical values are included in Table~\ref{tab:bo_details} once verified from the archived experiment configuration.}

\section{Experimental Design}

\subsection{Vehicle model, fleet, and tracking task}
\label{subsec:vehicle_task}
The study uses a fixed-speed lateral vehicle benchmark derived from the standard small-angle bicycle model~\cite{weber2026ifac}. The control-design model has state $\mathbf{x}=[\beta,r]^\top$, steering input $\delta$, and constant longitudinal speed $U$. Its continuous-time dynamics are
\begin{equation}
 \dot{\mathbf{x}}=\mathbf{A}\mathbf{x}+\mathbf{B}\delta,
 \label{eq:bicycle_ct}
\end{equation}
with
\begin{equation}
\mathbf{A}=\begin{bmatrix}
-\dfrac{C_f+C_r}{mU} & \dfrac{-aC_f+bC_r}{mU^2}-1\\[2mm]
\dfrac{-aC_f+bC_r}{I_z} & -\dfrac{a^2C_f+b^2C_r}{I_zU}
\end{bmatrix},
\label{eq:bicycle_A}
\end{equation}
and
\begin{equation}
\mathbf{B}=\begin{bmatrix}\dfrac{C_f}{mU}\\[2mm]\dfrac{aC_f}{I_z}\end{bmatrix}.
\label{eq:bicycle_B}
\end{equation}
The model is discretized by zero-order hold with sampling time $T_s=0.05$ s.

Controller candidates are evaluated on a nonlinear version of the same lateral plant. The front and rear tire slip angles are
\begin{align}
 \alpha_f&=\delta-\beta-\frac{ar}{U},\\
 \alpha_r&=-\beta+\frac{br}{U},
 \label{eq:slip_angles}
\end{align}
and each lateral force is computed as
\begin{equation}
 F_{y,q}=\mu F_{z,q}\tanh\left(\frac{C_q\alpha_q}{\mu F_{z,q}}\right),\qquad q\in\{f,r\}.
 \label{eq:tire_nonlinearity}
\end{equation}
\rev{With static axle loads $F_{z,f}=mgb/(a+b)$ and $F_{z,r}=mga/(a+b)$, the nonlinear states evolve according to}
\begin{align}
 \dot\beta&=\frac{F_{y,f}+F_{y,r}}{mU}-r,\\
 \dot r&=\frac{aF_{y,f}-bF_{y,r}}{I_z}.
 \label{eq:nonlinear_vehicle}
\end{align}
This makes the evaluation plant intentionally more nonlinear than the model used for LQI synthesis.

The fleet contains 30 systems in every realization, with ten systems from each of three families. Table~\ref{tab:vehicle_params} lists the family prototypes. Payload systems increase mass by 20\% and yaw inertia by 25\%. Tire-degraded systems reduce front cornering stiffness by 25\%. The generator additionally introduces small client-to-client variation around each family prototype using a variability setting of 0.025. Family labels are used only to construct and analyze the benchmark. They are not provided to BO.

\rev{The 0.025 variability setting perturbs the physical family prototype before simulation. The exact list of perturbed parameters, distribution, and any clipping or rejection rules are marked below for source-code verification because these details are not contained in the current paper-only branch.}
\begin{table}[t]
\centering
\caption{Vehicle-family prototypes and common simulation parameters.}
\label{tab:vehicle_params}
\scriptsize
\setlength{\tabcolsep}{3.5pt}
\begin{tabular}{@{}lccc@{}}
\toprule
Parameter & Nominal & Payload & Tire-degraded\\
\midrule
$m$ [kg] & 1500 & 1800 & 1500\\
$I_z$ [$\mathrm{kg\,m^2}$] & 2700 & 3375 & 2700\\
$a$ [m] & 1.25 & 1.25 & 1.25\\
$b$ [m] & 1.35 & 1.35 & 1.35\\
$C_f$ [$\mathrm{N/rad}$] & 80000 & 80000 & 60000\\
$C_r$ [$\mathrm{N/rad}$] & 90000 & 90000 & 90000\\
$U$ [m/s] & 15 & 15 & 15\\
$T_s$ [s] & 0.05 & 0.05 & 0.05\\
$\mu$ [-] & 1.0 & 1.0 & 1.0\\
\bottomrule
\end{tabular}
\end{table}
\rev{\textbf{Generator detail to verify before submission:} identify which of $m,I_z,a,b,C_f,C_r$ receive the 0.025 perturbation, the sampling distribution, and the positivity/rejection rule.}

\subsection{Identification data}
\rev{The identified model is obtained locally from the same fixed identification maneuver for every system, so the dynamics descriptor is not confounded by the two calibration maneuvers used later in BO. Recursive least squares estimates $\hat{\boldsymbol\theta}_i$ and its covariance from local identification data.}
\rev{\textbf{Identification detail to verify before submission:} add the exact parameter vector, identification-record length, excitation definition, RLS initialization, forgetting factor, and covariance convention from the archived experiment code or the corresponding setup in~\cite{weber2026ifac}.}

\subsection{Calibration task and feasibility}
Each BO evaluation measures yaw-rate root-mean-square error over a common pair of lane-change maneuvers. Their lane-change half-durations are $T_{\mathrm{lc}}=1.0$ s and $T_{\mathrm{lc}}=2.0$ s. Every controller candidate is evaluated on three stochastic episodes of each maneuver, giving six calibration episodes per BO query. All systems use this same balanced task distribution so that differences in the calibration objective arise from the plant rather than from different assigned maneuvers.

\begin{figure}[t]
\centering
\figurepdf{0.96\columnwidth}{figures/fig6_lane_change_reference.pdf}
\caption{Yaw-rate references used for controller calibration. Every system is evaluated on both lane-change maneuvers, and each BO query averages three stochastic episodes per maneuver.}
\label{fig:lane_change_reference}
\end{figure}

\rev{For a discrete trajectory with $T$ samples, the reported tracking metric is}
\begin{equation}
 J_i(\boldsymbol\xi,\omega)=\sqrt{\frac{1}{T}\sum_{k=1}^{T}\left(r_{i,k}(\boldsymbol\xi,\omega)-r_k^{\rm ref}\right)^2}.
 \label{eq:rmse}
\end{equation}
\rev{The two reference waveforms are generated by the same lane-change reference generator and differ only in $T_{\rm lc}$. Figure~\ref{fig:lane_change_reference} shows both signals. \textbf{Reference-waveform detail to verify before submission:} add the exact analytic or algorithmic construction and total episode duration.}

Yaw-rate RMSE is the calibration metric. A candidate is marked infeasible if LQI synthesis fails or if the nonlinear closed-loop rollout violates the prescribed limits. The sideslip magnitude is limited to $6^\circ$. The steering-rate change is limited to $0.2618$ rad per simulation step. These feasibility conditions are applied consistently to the baseline controller, source histories, offline reference search, and new-system evaluations.

\subsection{Landscape-similarity study}
\rev{To test whether dynamics heterogeneity translates into calibration heterogeneity, all 30 systems in each of five independent fleet realizations are evaluated on the same 128-point Sobol design $\{\boldsymbol\xi^{(m)}\}_{m=1}^{128}$. For system $i$, let $\mathcal F_i$ denote the indices of feasible calibration points and let $\mathcal F_{ij}=\mathcal F_i\cap\mathcal F_j$ denote the jointly feasible subset for a pair. Pairwise landscape similarity is computed only on $\mathcal F_{ij}$,}
\begin{equation}
 \rho_{ij}^{J}=\operatorname{Spearman}\left(\{J_i(\boldsymbol\xi^{(m)})\}_{m\in\mathcal F_{ij}},\{J_j(\boldsymbol\xi^{(m)})\}_{m\in\mathcal F_{ij}}\right).
 \label{eq:landscape_similarity}
\end{equation}
\rev{This construction prevents common infeasible points from artificially increasing the rank correlation. Feasibility-set agreement is analyzed separately with the Jaccard overlap $|\mathcal F_i\cap\mathcal F_j|/|\mathcal F_i\cup\mathcal F_j|$. The study contains 2175 within-realization system pairs. Realization-level summaries are computed within each of the five fleets and then bootstrapped across fleets so that the large number of dependent pairs is not treated as independent replication.}

\begin{figure*}[t]
\centering
\subfloat[Vehicle-parameter heterogeneity.\label{fig:landscape_params}]{\figurepdf{0.475\textwidth}{figures/fig2a_vehicle_parameter_heterogeneity.pdf}}
\hfill
\subfloat[Dynamics distance versus calibration-landscape similarity.\label{fig:landscape_similarity}]{\figurepdf{0.475\textwidth}{figures/fig2b_dynamics_vs_landscape_similarity.pdf}}
\caption{\rev{Plant heterogeneity produces only limited calibration-landscape heterogeneity. Panel (a) shows the physical variation across the three vehicle families. Panel (b) shows that rank correlation remains high over the full dynamics-distance range when computed only on jointly feasible calibration points.}}
\label{fig:landscape}
\end{figure*}

\subsection{Held-out fleet-informed BO evaluation}
The final experiment uses 20 fleet realizations that were not used for method development. Each realization contains 30 previously calibrated source systems, with ten systems from each family. Every source system contributes an independent BO history of 25 evaluations. Six systems per realization, two from each family, are then treated in turn as newly arriving systems. For each such system, its own completed history is removed and the remaining $29\times25=725$ source observations are available before its calibration begins. This produces 120 held-out new-system calibration problems.

Each new system starts with $n_{\rm init}=2$ local points from a Sobol low-discrepancy sequence. The system then receives at most ten BO-selected evaluations, giving a total local budget of 12 experiments. Independent BO, global warm start, and dynamics-weighted warm start use the same local initialization and the same evaluation randomness.

\rev{The experimental accounting is therefore intentionally asymmetric in available information but symmetric in newly executed target experiments. Independent BO receives $N_i^{\rm hist}=0$, while both warm-start methods receive $N_i^{\rm hist}=725$ records that were collected before the new system arrives. All three methods are allowed the same maximum $N_i^{\rm new}=12$ local experiments. This design evaluates incremental commissioning cost after prior fleet experience has already been accumulated.}

\rev{The primary metric is the budget-restricted $\Nfive$ defined above. As a secondary continuous metric, we use independently rescored reference-relative regret}
\begin{equation}
 r_i^{\mathrm{ref}}(n)=\frac{J_{i,\mathrm{test}}^{\mathrm{best}}(n)-J_i^{\mathrm{ref}}}{J_i^{\mathrm{ref}}}.
 \label{eq:regretref}
\end{equation}
\rev{We also report the area under this regret curve over the early local budget and fixed-budget probabilities of reaching the 5\% criterion. Fleet realization is the statistical replication unit. Metrics are first averaged over the six held-out systems in each realization, paired method differences are then formed for each of the 20 realizations, and 95\% confidence intervals are obtained by bootstrap resampling of these realization-level differences. The Wilcoxon signed-rank test is applied to the same 20 paired realization-level mean differences. \textbf{Bootstrap detail to verify before submission:} insert the exact number of bootstrap resamples used by the analysis script.}

\begin{table}[t]
\centering
\caption{Final held-out evaluation protocol.}
\label{tab:protocol}
\scriptsize
\begin{tabular}{@{}ll@{}}
\toprule
Quantity & Setting\\
\midrule
Source systems per realization & 30, 10 per family\\
Held-out fleet realizations & 20\\
New systems per realization & 6, 2 per family\\
Source history per system & 25 independent BO evaluations\\
$N_i^{\rm hist}$, warm start & 725 source observations\\
Local Sobol initialization & $n_{\rm init}=2$\\
Maximum $N_i^{\rm new}$ & 12 evaluations\\
Methods & independent / global / similarity\\
Primary metric & $\Nfive$\\
\bottomrule
\end{tabular}
\end{table}

\rev{A separate nine-source coverage experiment evaluates whether the benefit depends on having all 29 source systems available. That experiment varies the number of same-family sources from five down to zero while keeping the source-pool size fixed at nine. It is used only as a robustness analysis and not for tuning the final held-out comparison.}

\section{Results}

\subsection{Fleet warm start halves local experimentation}
\rev{Figure~\ref{fig:convergence} shows independently rescored reference-relative regret versus the number of new experiments performed on each held-out system. Both fleet-informed methods improve markedly faster than independent BO, while global and dynamics-weighted warm start remain nearly indistinguishable.}

\begin{figure*}[t]
\centering
\figurepdf{0.82\textwidth}{figures/fig3_confirmatory_convergence.pdf}
\caption{\rev{Held-out convergence over 20 fleet realizations and 120 new-system calibration problems after independently rescoring all incumbents on the reference-test bank. Fleet warm start accelerates controller calibration from the first BO iterations. Dynamics weighting and unweighted pooling exhibit essentially identical convergence.}}
\label{fig:convergence}
\end{figure*}

\rev{Under the strengthened reference protocol, independent BO has a budget-restricted mean time to the 5\% criterion of 6.717 local experiments, compared with 3.083 for global fleet warm start. This is a 54.1\% reduction in incremental commissioning experiments. The paired realization-level global-minus-independent difference is $-3.633$ experiments, with 95\% bootstrap CI $(-5.01,-2.37)$ and Wilcoxon $p=0.0005$. The favorable direction occurs in 15 of the 20 held-out fleet realizations and remains consistent across all three vehicle families. Independent BO reaches the criterion within the 12-evaluation budget for 95 of 120 targets, or 79.2\%, whereas both warm-start methods reach it for all 120 targets.}

\begin{table}[t]
\centering
\caption{\rev{Held-out new-system sample efficiency under the strengthened reference protocol.}}
\label{tab:main_results}
\scriptsize
\begin{tabular}{@{}lcc@{}}
\toprule
Method & Restricted mean $\Nfive$ & $P(T_{5\%}\le B)$\\
\midrule
Independent & 6.717 & 0.792\\
Global warm start & 3.083 & 1.000\\
Dynamics weighted & 3.083 & 1.000\\
\bottomrule
\end{tabular}
\end{table}

\rev{Figure~\ref{fig:n5} complements the restricted mean with the full target-attainment behavior. Panel (a) treats unresolved runs as right-censored, so the independent curve correctly plateaus below one instead of forcing censored runs to appear as successes at the final budget. Panel (b) reports success probabilities after three and five local experiments.}

\begin{figure*}[t]
\centering
\subfloat[Cumulative fraction reaching the 5\% criterion.\label{fig:n5_ecdf}]{\figurepdf{0.475\textwidth}{figures/fig4a_n5pct_ecdf.pdf}}
\hfill
\subfloat[Success probability at fixed local budgets.\label{fig:n5_success}]{\figurepdf{0.475\textwidth}{figures/fig4b_n5pct_success_probability.pdf}}
\caption{\rev{New-system experiments required to reach within 5\% of the independently scored system-specific offline reference. Panel (a) shows cumulative observed target attainment under the common 12-evaluation censoring budget. Panel (b) shows the probability of meeting the criterion within three and five local experiments.}}
\label{fig:n5}
\end{figure*}

\subsection{Dynamics-aware weighting adds no measurable benefit}
\rev{On the held-out fleets, dynamics-weighted warm start provides no detectable improvement over global pooling. The paired realization-level mean difference in budget-restricted $\Nfive$ is $0.000$ experiments with 95\% CI $(-0.05,+0.04)$ and Wilcoxon $p=1.000$. The independently rescored convergence curves in Fig.~\ref{fig:convergence} likewise remain nearly coincident.}

\rev{The mechanistic study explains why this null result is plausible without implying that dynamics information is completely irrelevant. On the 128-point jointly feasible landscape grid, 99.4\% of all 2175 system pairs satisfy $\rho_{ij}^{J}\ge0.95$, with a minimum of 0.9355. Mean rank correlation is 0.991 for same-family pairs and 0.985 for cross-family pairs, while the median feasibility-set Jaccard similarity is 1.0. Dynamics similarity $s_{ij}$ still has a modest positive association with landscape similarity. The median within-realization Spearman correlation is 0.245 with 95\% CI $(0.162,0.288)$. Thus the calibration landscapes are not exactly invariant, but the residual dependence is small relative to their uniformly high common correlation floor. This is consistent with the repeated observation that dynamics-aware weighting does not improve outcome-level sample efficiency over broad pooling.}

\begin{figure*}[t]
\centering
\subfloat[Held-out global versus dynamics-weighted warm start.\label{fig:null_pairwise}]{\figurepdf{0.475\textwidth}{figures/fig5a_global_vs_weighted_parity.pdf}}
\hfill
\subfloat[Sensitivity to same-family source coverage.\label{fig:null_coverage}]{\figurepdf{0.475\textwidth}{figures/fig5b_coverage_sensitivity.pdf}}
\caption{Dynamics-aware source weighting does not improve new-system sample efficiency. Panel (a) compares global and weighted warm start on the held-out fleets. Panel (b) shows that their performance remains similar when same-family source coverage is deliberately reduced.}
\label{fig:null_weighting}
\end{figure*}

\subsection{Continuous early-budget performance}
\rev{The 5\% threshold is intentionally easy to interpret, but it still selects one performance level. We therefore use the area under the independently rescored reference-relative regret curve as a secondary continuous measure of early calibration progress. Figure~\ref{fig:convergence} shows the same ordering as the threshold metric: both warm-start methods reduce regret substantially earlier than independent BO, while global and dynamics-weighted transfer remain nearly indistinguishable. The agreement between the threshold-based and continuous views shows that the sample-efficiency conclusion is not an artifact of the 5\% criterion.}

\section{Discussion}
\rev{The main practical conclusion is that an existing fleet should be treated as a reusable source of controller-calibration information. In the studied LQI problem, a newly arriving system reaches near-reference performance after roughly three local experiments when historical fleet observations are available, compared with a budget-restricted mean of 6.7 experiments under independent BO. All 120 warm-started target problems reach the 5\% criterion within the 12-evaluation budget, whereas 20.8\% of the independent runs remain unresolved at that point.}

\rev{The more surprising result concerns heterogeneity. The vehicle families differ in mass, yaw inertia, and tire properties, and Fig.~\ref{fig:landscape} shows that their identified dynamics are clearly distinguishable. Nevertheless, these differences translate into only limited changes in the ordering of LQI calibration settings. On the denser jointly feasible landscape study, 99.4\% of system pairs remain above $\rho_{ij}^{J}=0.95$. A modest residual relationship with dynamics similarity remains, so plant and calibration heterogeneity are not completely decoupled. For the fixed-speed, three-parameter LQI calibration problem and maneuver family studied here, however, that residual dependence is too small to yield a measurable advantage for dynamics-weighted transfer. This interpretation should not be generalized to higher-dimensional controller classes, different operating conditions, or arbitrary heterogeneous plants without further evidence.}

\rev{This observation also explains why simple pooling performs so well. In a three-dimensional search space with little new-system data, additional historical observations improve GP coverage even when they originate from dynamically different systems. The result should not be read as evidence that dynamics similarity contains no information. Rather, the remaining calibration-landscape differences are small relative to the strong cross-system agreement, so selectively attenuating dynamically dissimilar records does not improve commissioning efficiency in the investigated benchmark.}

The historical source records are not free. Each was generated by experiments performed when earlier systems were commissioned. The economic argument is instead one of reuse. Once these evaluations have been collected, their cost can be distributed across subsequent systems. \rev{Accordingly, the reported 54\% reduction is a reduction in incremental new-system commissioning experiments, not a 54\% reduction in total fleet-wide data generation.}

\rev{Several limitations remain. Validation is simulation based. The nonlinear evaluation plant is intentionally mismatched to the controller-design model, but real-vehicle confirmation is still missing. The study considers one controller structure, one three-dimensional calibration parameterization, one operating speed, and one lateral tracking application. Higher-dimensional or more task-specific calibration problems may exhibit less transferable landscapes. The dense offline reference is numerically well converged and independently scored, but it remains a numerical reference rather than a proof of the global optimum. The present architecture assumes that scalar historical calibration records can be stored centrally or otherwise made available to the new system and does not provide a formal privacy guarantee.}

\section{Conclusion}
\rev{This paper investigated whether controller-calibration experience accumulated by previously commissioned fleet members can reduce the commissioning effort for a newly arriving system. In a held-out study on nonlinear vehicle lateral dynamics, warm-starting Bayesian optimization with historical fleet observations reduced the budget-restricted mean number of new-system experiments required to reach within 5\% of an independently scored system-specific offline reference from 6.72 to 3.08, while all warm-started systems reached the criterion within the 12-evaluation budget. Dynamics-aware source weighting provided no measurable improvement over simple global pooling. A denser calibration-landscape analysis explains this result more precisely: 99.4\% of system pairs remain rank-correlated above 0.95 on jointly feasible calibration points despite substantial dynamics heterogeneity, although a modest residual relationship with dynamics similarity persists. The result therefore supports broad reuse of historical fleet experience for this calibration problem while showing that the heterogeneity relevant for plant modeling need not translate into equally strong heterogeneity in controller-parameter optimization.}

\bibliographystyle{IEEEtran}
\bibliography{references}

\end{document}
